\SourcePage{474}
\SGAsection{12}{Jacobian 的 Néron 模型与 Picard--Lefschetz 公式}

先回顾 M. Raynaud 的下列\Term{generic}{一般}结果 [22，\Term{theorem}{定理} 5 与\Term{theorem}{定理} 7；22 bis]。

\begin{SGAresult}{\Term{theorem}{定理}}{12.1（M. Raynaud）}
设 $S$ 是\Term{strictly-local}{严格局部}的\Term{trait}{离散赋值环谱}，$f:X\to S$ 是 $S$ 上的\Term{proper}{固有}\Term{flat}{平坦}\Term{scheme}{概形}，其\Term{fibre}{纤维}\Term{dimension}{维数}为 $1$，满足
\[
 f_*(\mathcal O_X)\simeq\mathcal O_S,\tag{12.1.1}
\]
并且 $X$ 的\Term{local-ring}{局部环}都是\Term{factorial}{析因环}。对 $X_o$ 的每个\Term{maximal-point}{极大点} $x_i$，令
\[
 d_i=\operatorname{length}\mathcal O_{X_o,x_i},\qquad
 d_i^t=\mu_i d_i,\tag{12.1.2}
\]
其中 $\mu_i$ 是 $k(x_i)$ 在 $k$ 上的\Term{radicial-multiplicity}{根性重数}（EGA IV 4.7.4），它是 $k$ 的\Term{characteristic-exponent}{特征指数}的一个幂。令 $d$（相应地 $d^t$）为整数 $d_i$（相应地 $d_i^t$）的最大公因数；因此，若 $k$ 是\Term{perfect-field}{完美域}，则 $d=d^t$。置
\[
 P=\underline{\operatorname{Pic}}_{X/S}:(\mathrm{Sch})_{/S}^{\mathrm o}\longrightarrow(\mathrm{Ab}).\tag{12.1.3}
\]
\textup{a)} 考虑下列条件：

1) $d^t=1$。

2) $X$ 有一个\Term{degree}{次数}为 $1$ 的\Term{zero-cycle}{$0$ 维圈}。

3) $X\times P^0$ 上存在一个\Term{invertible-module}{可逆模} $L$，使对应的\Term{homomorphism}{同态}
$P_\eta^0\to\underline{\operatorname{Pic}}_{X_\eta/\eta}=P_\eta$ 是\Term{canonical}{典范}\Term{inclusion}{包含}。（注意，已知 $P_\eta$ 可由 $\eta$ 上的\Term{locally-finite-type}{局部有限型}\Term{group-scheme}{群概形}\Term{representable}{表示}；这是 J. P. Murre 的结果，例如参见 SGA 6 第 XII 讲 1.5。）

\SourcePage{475}
4) $X$ 在 $S$ 上\Term{cohomologically-flat}{上同调平坦}（EGA III 7.8.1；也就是说，在\Term{dimension}{维数} $0$ 上\Term{cohomologically-flat}{上同调平坦}，或者说 $k\xrightarrow\sim H^0(X_o,\mathcal O_{X_o})$），并且 $d=1$。

5) $P$ \Term{representable}{可表}（或者说，可由 $S$ 上的\Term{smooth}{光滑}\Term{group-scheme}{群概形}\Term{representable}{表示}）。

6) $P^0$ \Term{representable}{可表}（或者说，可由 $S$ 上的\Term{smooth}{光滑}\Term{group-scheme}{群概形}\Term{representable}{表示}）。（关于 $P^0$ 的定义，参见 SGA 3 第 VI$_{\mathrm B}$ 讲 3.1。）

7) $P^0$ 在 $S$ 上\Term{separated}{分离}，即满足 EGA II 7.2.3 的\Term{valuative-criterion-separatedness}{分离性的赋值判别准则}（或者说，$P^0$ 在 $S$ 上满足 $s(\eta)=0$ 的任一\Term{section}{截面} $s$ 都为零）。

8) $d=1$。

则有如下蕴含关系：
\[
 1)\Longleftrightarrow2)\Longrightarrow3)\Longrightarrow4)
 \Longleftrightarrow5)\Longleftrightarrow6)\Longleftrightarrow7)\Longrightarrow8).
 \tag{12.1.4}
\]
特别地，若 $d=d^t$，例如 $k$ 是\Term{perfect-field}{完美域}，则条件 1) 至 8) \Term{equivalence}{等价}。

\textup{b)} 假设 $d^t=1$，因而 $P$ 可由 $S$ 上的\Term{smooth}{光滑}\Term{group-scheme}{群概形}\Term{representable}{表示}。则 $P$ 的\Term{unit-section}{单位截面}的\Term{schematic-closure}{概形闭包}是 $S$ 上的一个\Term{etale}{平展}\Term{subgroup-scheme}{子群概形} $E$。若 $I$ 是 $X_o$ 的\Term{reduced}{约化}\Term{irreducible-component}{不可约分支} $X_o^i$ 的\Term{set}{集合}，则得到\Term{isomorphism}{同构}
\[
 \mathbb Z^I\xrightarrow\sim\Gamma(S,E),\tag{12.1.5}
\]
\SourcePage{476}
它将每个 $i\in I$ 对应于由 $S$ 上的\Term{invertible-sheaf}{可逆层} $\mathcal O_X(X_o^i)$ 定义的 $P$ 在 $S$ 上的\Term{section}{截面} $s_i$，此\Term{section}{截面}在 $\eta$ 处为零。（注意，按假设 $X$ 的\Term{local-ring}{局部环}都是\Term{factorial}{析因环}，故\Term{cycle}{圈} $X_o^i$ 确实是\Term{divisor}{除子}；因此，可以将 $\mathcal O_X(X_o^i)$ 理解为这个\Term{divisor}{除子}对应的\Term{invertible-sheaf}{可逆层}，见 EGA IV 21.2.8.1。）并且有
\[
 P^0\cap E=P\text{ 的零截面}.\tag{12.1.6}
\]
\textup{c)} 在 b) 的条件下，\Term{quotient}{商}
\[
 Q=P/E\tag{12.1.7}
\]
可由 $S$ 上的\Term{smooth}{光滑}\Term{separated}{分离}\Term{group-scheme}{群概形}\Term{representable}{表示}，而且后者具有\Term{neronian}{Néron 性}，即满足用其\Term{special-fiber}{特殊纤维}表述的\Term{universal-property}{泛性质} (1.1.2)。

\textup{d)} 在 c) 的条件下，进一步假设 $X$ 的\Term{geometric-generic-fiber}{几何泛纤维}是\Term{integral}{整的}。则有\Term{canonical}{典范}\Term{exact-sequence}{正合列}
\[
 0\longrightarrow A\longrightarrow Q\xrightarrow{\deg}\mathbb Z_S\longrightarrow0,\tag{12.1.8}
\]
其中 $\deg$ 由“\Term{degree}{次数}”\Term{homomorphism}{同态} $P\to\mathbb Z_S$ 诱导，$A$ 是 $S$ 上的\Term{smooth}{光滑}\Term{separated}{分离}\Term{finite-type}{有限型}\Term{group-scheme}{群概形}，也具有\Term{neronian}{Néron 性}，即它是
\[
 A_\eta=P_\eta^0=\underline{\operatorname{Pic}}_{X_\eta/\eta}^0\tag{12.1.9}
\]
的\Term{neron-model}{Néron 模型}。另一方面，复合\Term{homomorphism}{同态} $P^0\to P\to Q$ 诱导一个\Term{isomorphism}{同构}
\[
 P^0\xrightarrow\sim A^0.\tag{12.1.10}
\]
\end{SGAresult}

\SourcePage{477}
\textbf{12.1.11。}我们关心的情形是：$X_\eta$ 在 $\eta$ 上\Term{smooth}{光滑}，从而由 (12.1.1)，$X$ 的\Term{geometric-generic-fiber}{几何泛纤维}\Term{smooth}{光滑}\Term{connected}{连通}，因而是\Term{integral}{整的}；并且 $X_\eta$ 有一个\Term{degree}{次数}为 $1$ 的\Term{zero-cycle}{$0$ 维圈}。这时，12.1 中考虑的全部条件，特别是 a) 中的条件 1) 至 8)，均成立。此外，$A_\eta=P_\eta^0$ 是\Term{abelian-scheme}{阿贝尔概形}，d) 给出了它的\Term{neron-model}{Néron 模型} $A$ 的十分明确的构造。特别地，(12.1.10) 给出\Term{isomorphism}{同构}
\[
 A_o^0\xrightarrow\sim\underline{\operatorname{Pic}}_{X_o/k}^0,\tag{12.1.12}
\]
这尤其表明：$A$ 在 $S$ 上具有\Term{semistable-reduction}{半稳定约化}，当且仅当 $\underline{\operatorname{Pic}}_{X_o/k}^0$ 的\Term{unipotent-rank}{幺幂秩}为零。当 $X$ \Term{regular}{正则}且 $X_o$ 在 $k$ 上\Term{separable}{可分}时，已知这又\Term{equivalence}{等价}于\Term{geometric-special-fiber}{几何特殊纤维} $X_{\bar o}$ 的\Term{singular-point}{奇点}全为\Term{ordinary-quadratic-point}{普通二次点}，即它们的\Term{completed-local-ring}{完备局部环}与\Term{affine-space}{仿射空间} $\mathbb E^2$ 的两条坐标轴之并在原点的\Term{completed-local-ring}{完备局部环}相同。利用 EGA IV 21.8.5，结合域 $k$ 上的\Term{proper}{固有}\Term{separable}{可分}\Term{curve}{曲线}的\Term{picard-scheme}{Picard 概形}的“众所周知”的描述，可以验证这一事实；将此留作读者的练习。另外还应注意，由\Term{semistable-reduction}{半稳定约化}的 Galois 判别准则 3.5，“$X_{\bar o}$ 的\Term{singular-point}{奇点}为\Term{ordinary-quadratic-singularity}{普通二次奇点}”蕴含“$A_\eta=\underline{\operatorname{Pic}}_{X_\eta/\eta}^0$ 具有\Term{semistable-reduction}{半稳定约化}”，也是“几何\Term{monodromy}{单值性}\Term{theorem}{定理}”（第 III 讲\SGAsourcegap{原 PDF 第477页此处在“III”后留空，缺少条款号}）的直接推论；但这个论证不能给出重要的\Term{isomorphism}{同构} (12.1.10) 及其推论 (12.1.12)。

\SourcePage{478}
\textbf{12.2。}关于 12.1 的证明，参见 [22 bis]；我们实际上也只需要 $X_\eta$ 在 $\eta$ 上\Term{smooth}{光滑}的情形。顺便指出，[22 bis] 在远比 12.1 c)、d) 宽泛的条件下证明了\Term{picard-scheme}{Picard 概形}与\Term{neron-model}{Néron 模型}之间的关系；由于其表述更为技术化，这里不将它们列入回顾。只指出，b)、c)、d) 都是 a) 的容易推论，而 a) 中只有蕴含关系 $3)\Rightarrow4)\Rightarrow5)$ 以及 $7)\Rightarrow4)$ 较为微妙，其中第一个最为困难。不过，若假设 $X_o$ 在 $k$ 上\Term{separable}{可分}，这一点便是显然的（因为 $S$ \Term{strictly-local}{严格局部}，这时 $X$ 在 $S$ 上有\Term{section}{截面}，故条件 1)，即 $d^t=1$，成立）。建议读者在这一特别简单的情形给出证明；必要时还可假设 $X_\eta$ \Term{smooth}{光滑}，因而 $P_\eta^0$ 是\Term{abelian-scheme}{阿贝尔概形}，这样就可以利用后者的\Term{neron-model}{Néron 模型} $A$，并以 $A^0$ \Term{representable}{表示} $\underline{\operatorname{Pic}}_{X/S}^0$。实际上，我们只在这一特殊情形使用 12.1，并且还假设 $X_{\bar o}$ 的\Term{singular-point}{奇点}只有\Term{ordinary-quadratic-point}{普通二次点}，即由 (12.1.11)，$A_\eta$ 具有\Term{semistable-reduction}{半稳定约化}。这就是 12.5.4 中考虑的情形，从现在起我们只讨论这一情形。

\textbf{12.3。}设 $C$ 是域 $k$ 上的\Term{separable}{可分}\Term{proper}{固有}\Term{curve}{曲线}。为简便起见，假设 $k$ \Term{separably-closed}{可分闭}，读者若愿意，也可以假设它\Term{algebraically-closed}{代数闭}。设 $\widetilde C$ 是 $C$ 的\Term{normalization}{正规化}。已知
\SourcePage{479}
$\underline{\operatorname{Pic}}_{\widetilde C/k}^0$ 的“\Term{reductive-rank}{约化秩}”为零，即其\Term{maximal-torus}{极大环面}为零。因此，$\underline{\operatorname{Pic}}_{C/k}^0$ 的\Term{maximal-torus}{极大环面}包含在
\[
 N=\operatorname{Ker}\bigl(\underline{\operatorname{Pic}}_{C/k}\longrightarrow
 \underline{\operatorname{Pic}}_{\widetilde C/k}\bigr)
\]
中，故与后者的\Term{maximal-torus}{极大环面}\Term{canonical}{典范}\Term{isomorphism}{同构}。而借助 EGA IV 21.8.5，很容易明确写出 $N$ 的结构：$N$ \Term{canonical}{典范}地是一个\Term{torus}{环面} $T$ 被一个\Term{connected}{连通}\Term{unipotent}{幺幂群} $U$ 的\Term{extension}{扩张}，其中 $T$ 的描述将在下文给出；$U$ 同样在 $k$ 上有一个\Term{composition-series}{合成列}，各商\Term{isomorphism}{同构}于\Term{additive-group}{加法群} $\mathbb G_a$。若 $T'$ 是 $N$ 的\Term{maximal-torus}{极大环面}，则 $T'\cap U$ 既\Term{unipotent}{幺幂}又是\Term{multiplicative-type}{乘法型}，因而为零（SGA 3 第 XVII 讲 2.4）；而 $T'$ 在 $T$ 中的像就是 $T$（SGA 3 第 XII 讲 7.1 e)），所以 $T'\to T$ 是\Term{isomorphism}{同构}。因此，$T$ \Term{canonical}{典范}\Term{isomorphism}{同构}于 $N$ 的\Term{maximal-torus}{极大环面}，也就\Term{canonical}{典范}\Term{isomorphism}{同构}于 $\underline{\operatorname{Pic}}_{X/k}^0$ 的\Term{maximal-torus}{极大环面}。

为了明确描述 $N$ 的结构，假设 $k$ \Term{separably-closed}{可分闭}会较为方便；此时 $C$ 的\Term{irreducible-component}{不可约分支} $C_i$ 都\Term{geometrically-irreducible}{几何不可约}，且 $C$ 的非\Term{normal}{正规}点在 $k$ 上都是\Term{radicial}{根性的}。令 $J$ 为 $\widetilde C$ 的\Term{irreducible-component}{不可约分支}的\Term{set}{集合}；对每个 $x\in C$，令 $J(x)$ 为 $C$ 经过 $x$ 的“\Term{branch}{分支}”的\Term{set}{集合}，即 $\widetilde C$ 中位于 $x$ 上方的点的\Term{set}{集合}；最后，令
\[
 R(x)=\mathbb Z^{J(x)}/\operatorname{Im}\mathbb Z,\tag{12.3.1}
\]
其中 $\mathbb Z\to\mathbb Z^{J(x)}$ 是\Term{diagonal-map}{对角映射}，并令
\[
 R=\left(\bigoplus_x R_x\right)/\operatorname{Im}\mathbb Z^J.\tag{12.3.2}
\]
\SourcePage{480}
这里显然只需对 $C$ 中同时属于至少两个\Term{branch}{分支}的\Term{finite}{有限}多个点求和。因此 $R$ 是一个\Term{finite-type}{有限型} $\mathbb Z$-\Term{module}{模}。于是有\Term{canonical}{典范}\Term{isomorphism}{同构}
\[
 T\ (\simeq\underline{\operatorname{Pic}}_{C/k}^0\text{ 的极大环面})
 \simeq R\otimes\mathbb G_{mk}.\tag{12.2.3}
\]
换句话说，$T$ 的\Term{character-group}{特征标群} $M$ 由下式给出：
\[
 M\overset{\mathrm{def}}=D(T)=R^\vee
 \simeq\operatorname{Ker}\left(\bigoplus_x R(x)^\vee\longrightarrow\mathbb Z^J\right),\tag{12.2.4}
\]
其中 $R(x)^\vee$ 也可理解为 $J(x)$ 上所有取整数值且各值之和为零的函数所成的\Term{set}{集合}。

\textbf{12.3.5。}我们只需要使用 $C$ 的每个点 $x$ 至多有两个\Term{branch}{分支}经过的情形，\Term{ordinary-quadratic-point}{普通二次点}尤其如此！对满足 $\operatorname{card}J(x)=2$ 的点，$R(x)$ 因而\Term{isomorphism}{同构}于 $\mathbb Z$；但这个\Term{isomorphism}{同构}依赖于为经过 $x$ 的两个\Term{branch}{分支} $(C',C'')$ 选择一个次序，\Term{equivalence}{等价}地，选择其中一个\Term{branch}{分支} $C'$（将它对应于 $0_{C'}+1_{C''}$ 在 $R(x)$ 中的\Term{canonical}{典范}像）。于是得到\Term{isomorphism}{同构}
\[
 R=\mathbb Z^I/\operatorname{Im}\mathbb Z^J,\quad
 T\simeq\mathbb G_m^I/\operatorname{Im}\mathbb G_m^J,\quad
 M=\operatorname{Ker}(\mathbb Z^I\longrightarrow\mathbb Z^J),\tag{12.3.6}
\]
其中 $I$ 是 $C$ 中属于两个\Term{branch}{分支}的点 $x$ 的\Term{set}{集合}，并且须如上所述选择\Term{signature}{符号}。

\textbf{注记 12.3.7。}有时，以下方式很便于明确描述 $\underline{\operatorname{Pic}}_{C/k}^0$ 的\Term{maximal-torus}{极大环面} $T$ 的\Term{character-group}{特征标群} $M$。考虑“一维（无向）\Term{graph}{图}”（N. Bourbaki，《群与 \Term{lie-algebra}{Lie 代数}》，第 IV 章，附录）
\SourcePage{481}
$\Gamma(C)$：其顶点是 $\widetilde C$ 的\Term{irreducible-component}{不可约分支}；两个顶点 $\widetilde C_i$、$\widetilde C_j$ 之间的边与\Term{set}{集合} $\widetilde C_i\mathbin{\times_C'}\widetilde C_j$ 一一对应，这个\Term{set}{集合}由 $C$ 中位于同一点、分别包含在 $\widetilde C_i$、$\widetilde C_j$ 内的两个不同\Term{branch}{分支}组成。当 $i\ne j$ 时，它就是 $\widetilde C_i\times_C\widetilde C_j$；当 $i=j$ 时，它是
$\widetilde C_i\times_C\widetilde C_i-\Delta(\widetilde C_i)$ 关于\Term{symmetric}{对称}关系的\Term{quotient}{商}，其中 $\Delta:\widetilde C\to\widetilde C\times_C\widetilde C$ 是\Term{diagonal-map}{对角映射}（因此，若 $C_i$ \Term{geometrically-unibranch}{几何单分支}，这个\Term{set}{集合}为空）。于是有\Term{canonical}{典范}\Term{isomorphism}{同构}
\[
 M=D(T)\xleftarrow\sim H_1(\Gamma(C),\mathbb Z),\qquad
 M^\vee=R\xrightarrow\sim H^1(\Gamma(C),\mathbb Z).\tag{12.3.7}
\]
我们不需要 $M$ 的这一组合解释，只定义\Term{canonical}{典范}\Term{homomorphism}{同态}
\[
 H_1(\Gamma(C),\mathbb Z)\longrightarrow M,
\]
并将它是\Term{isomorphism}{同构}的证明留给读者。令 $I$ 为 $C$ 中至少属于两个\Term{branch}{分支}的点的\Term{set}{集合}，$J$ 如前，为 $\widetilde C$ 的\Term{irreducible-component}{不可约分支}的\Term{set}{集合}（也可取 $C$ 的\Term{irreducible-component}{不可约分支}，二者相同），$K$ 为 $C$ 在 $I$ 中各点的“\Term{branch}{分支}”的\Term{set}{集合}，即 $I$ 在 $\widetilde C$ 中的\Term{inverse-image}{逆像}。于是有\Term{canonical-map}{典范映射}
\[
 h:K\longrightarrow I\times J,\qquad h=(f,g),\tag{12.3.9}
\]
而 (12.3.4) 中对 $M$ 的描述可用 $h$ 表为
\[
 M\simeq\operatorname{Ker}\bigl(\mathbb Z^{(h)}:
 \mathbb Z^{(K)}\longrightarrow\mathbb Z^{(I\times J)}\bigr),\tag{12.3.10}
\]
因为 (12.3.3) 中 $R=M^\vee$ 的表达式显然可以理解为
\SourcePage{482}
由 $h$ 诱导的\Term{homomorphism}{同态} $\mathbb Z^{I\times J}\to\mathbb Z^K$ 的\Term{cokernel}{余核}，此\Term{homomorphism}{同态}是 $\mathbb Z^{(h)}$ 的\Term{transpose}{转置}。还应注意，\Term{graph}{图} $\Gamma(C)$ 也可以显然地用 $h$ 描述：
\[
 \Gamma_0=J,\qquad
 \Gamma_1=(K\times_I K-\Delta(K))/\text{对称关系}.
\]
另一方面，为 $\Gamma(C)$ 的每条边选择一个\Term{orientation}{定向}，并记这些边 $(x_i,C',C'')$ 的\Term{set}{集合}为 $\Gamma_1$，则
\[
 H_1(\Gamma(C),\mathbb Z)
 =\operatorname{Ker}\bigl(d:\mathbb Z^{(\Gamma_1)}\longrightarrow\mathbb Z^{(J)}\bigr),\tag{12.3.11}
\]
其中\Term{differential-operator}{微分算子}由下式给出：
\[
 d(x,C',C'')=g(C'')-g(C').\tag{12.3.12}
\]
于是，\Term{homomorphism}{同态} (12.3.8) 由\Term{chain-complex}{链复形}的\Term{homomorphism}{同态}
\[
 \begin{split}
 v=(v_0,v_1):C_\bullet(\Gamma(C))
 &=[\mathbb Z^{(\Gamma_1)}\longrightarrow\mathbb Z^{(J)}]\\
 &\longrightarrow[\mathbb Z^{(K)}\longrightarrow\mathbb Z^{(I\times J)}]
 \end{split}\tag{12.3.13}
\]
诱导，其中
\[
 v_1(x,C',C'')=(x,C'')-(x,C'),\qquad
 v_0(\widetilde C_j)=(0,\widetilde C_j).\tag{12.3.14}
\]

\textbf{12.4。}现在置于 12.2 末尾考虑的情形。由 (12.1.12)，并将 (12.3.6) 用于 $C=X_o$，$A_\eta=\underline{\operatorname{Pic}}_{X_\eta/\eta}^0$ 的\Term{neron-model}{Néron 模型} $A$ 的\Term{special-fiber}{特殊纤维} $A_o^0$ 的\Term{maximal-torus}{极大环面} $T_o$ 的\Term{character-group}{特征标群} $M$ 由下式给出：
\[
 M=\operatorname{Ker}(\mathbb Z^I\xrightarrow d\mathbb Z^J),\tag{12.4.1}
\]
其中 $I$ 是 $X_o$ 的\Term{singular-point}{奇点}\Term{set}{集合}，$J$ 是其\Term{irreducible-component}{不可约分支}\Term{set}{集合}，\Term{homomorphism}{同态} $\mathbb Z^I\to\mathbb Z^J$ 由
\SourcePage{483}
\[
 d(x_i)=C_i'-C_i''\tag{12.4.2}
\]
给出。这里 $C_i'$、$C_i''$ 是 $\widetilde C$ 中对应于经过\Term{double-point}{二重点} $x_i$ 的两个\Term{branch}{分支}的\Term{irreducible-component}{不可约分支}。\Term{operator}{算子} $d$ 依赖于为这两个\Term{branch}{分支}选择的次序 $(C_i',C_i'')$，\Term{isomorphism}{同构} (12.4.1) 显然也依赖于这些选择。

不过，若对 $i\in I$，记
\[
 \varphi_{x_i}:M\longrightarrow\mathbb Z\tag{12.4.3}
\]
为复合映射 $M\to\mathbb Z^I\xrightarrow{\operatorname{pr}_i}\mathbb Z$，其中第一个\Term{homomorphism}{同态}来自 (12.4.1)，则回到表达式 (12.3.4) 便立即看出，$\varphi_{x_i}$ 在相差一个\Term{signature}{符号}的意义下完全确定；更确切地说，它只依赖于为经过 $i$ 的两个\Term{branch}{分支} $C_i'$、$C_i''$ 选择的次序，而交换次序时它也变号。因此，特别地，\Term{bilinear-form}{双线性型}
\[
 \varphi_{x_i}\otimes\varphi_{x_i}
 \in M^\vee\otimes_{\mathbb Z}M^\vee
 \simeq\operatorname{Hom}(M\otimes M,\mathbb Z)\tag{12.4.4}
\]
由\Term{singular-point}{奇点} $x_i$ \Term{canonical}{典范}地确定，不再有\Term{signature}{符号}的歧义。于是可以置
\[
 u=\sum_i\varphi_{x_i}\otimes\varphi_{x_i}:
 M\otimes_{\mathbb Z}M\longrightarrow\mathbb Z,\tag{12.4.5}
\]
从而得到一个完全确定、即不依赖于任何\Term{signature}{符号}选择的\Term{pairing}{配对}。它显然\Term{positive-definite}{正定}，因为它是 $\mathbb Z^I$ 上的标准\Term{quadratic-form}{二次型} $\sum_i X_i^2$ 在 $M\hookrightarrow\mathbb Z^I$ 上的\Term{restriction}{限制}。

\SourcePage{484}
将 $T_o$ 的\Term{character-group}{特征标群}记为 $M$，乍看与前面各节的记号不一致：那里将它记为 $M'$，而以 $M$ 表示与 $A_K$ 的\Term{dual}{对偶}\Term{abelian-scheme}{阿贝尔概形} $A_K'$ 相关的相应\Term{group}{群}。但 $A_K$ 是一个\Term{jacobian}{Jacobian}，因此带有一个\Term{canonical}{典范}\Term{polarization}{极化} $\xi$ [5]，诱导\Term{canonical}{典范}\Term{isomorphism}{同构}
\[
 \psi_\xi:A_K\xrightarrow\sim A_K'.\tag{12.4.6}
\]
利用这个\Term{isomorphism}{同构}，将 $A_K$ 与 $A_K'$ 认同，就无须再区分 $M$ 与 $M'$。特别地，对任一\Term{prime-number}{素数} $\ell$，\Term{monodromy-pairing}{单值性配对} (9.1.2) 可视为\Term{pairing}{配对}
\[
 u_\ell:M_\ell\otimes_{\mathbb Z_\ell}M_\ell\longrightarrow\mathbb Z_\ell.\tag{12.4.7}
\]
于是可以陈述本节的主要结果。

\begin{SGAresult}{\Term{theorem}{定理}}{12.5（\Term{picard-lefschetz-formula}{Picard--Lefschetz 公式}）}
设 $S$ 是\Term{strictly-local}{严格局部}的\Term{trait}{离散赋值环谱}，$X$ 是 $S$ 上的\Term{regular}{正则}\Term{projective}{射影}\Term{flat}{平坦}\Term{scheme}{概形}，其\Term{fibre}{纤维}\Term{dimension}{维数}为 $1$，\Term{geometric-generic-fiber}{几何泛纤维}\Term{smooth}{光滑}\Term{connected}{连通}，\Term{geometric-special-fiber}{几何特殊纤维}的\Term{singular-point}{奇点}全为\Term{ordinary-quadratic-point}{普通二次点}。沿用 12.4 引入的记号，对任一\Term{prime-number}{素数} $\ell$，\Term{monodromy-pairing}{单值性配对} (12.4.7) 等于从\Term{pairing}{配对} (12.4.5) 经 $\mathbb Z\to\mathbb Z_\ell$ \Term{extension-of-scalars}{扩张标量}得到的\Term{pairing}{配对}。
\end{SGAresult}

\SourcePage{485}
\textbf{12.6。}注意，在所考虑的特殊情形，若在 10.4 b) 中取 $\xi$ 为 10.4 中考虑的\Term{canonical}{典范}\Term{polarization}{极化}，这个陈述就证明了 10.4。如 10.5 所示，这已足以\Term{pushout}{推出}\Term{generic}{一般}情形的 10.4。

\textbf{12.7。12.5 的证明：$\ell\ne p$ 的情形。}证明主要依赖于 12.5 的情形下的\Term{cohomological-picard-lefschetz}{上同调 Picard--Lefschetz 公式}。当 $\ell\ne p$ 时，该公式给出\Term{monodromy-group}{单值群} $I$ 在
\[
 H^1(X_{\bar\eta},\mathbb Z_\ell)
 \simeq H^1(A_{\bar\eta},\mathbb Z_\ell)
 \simeq\operatorname{Hom}(T_\ell(A_{\bar\eta}),\mathbb Z_\ell)\tag{12.7.1}
\]
上的\Term{action}{作用}。事实上，由 9.2 中 $\ell\ne p$ 时 $u_\ell$ 的定义，该公式对所考虑的 $\ell$ 而言，与断言 12.5 实质相同；它将在后面的一讲（第 XV 讲）中，在实际上宽泛得多的情形下证明。于是只剩 $\ell=p$ 的情形。我们将用归约到\Term{characteristic-zero}{特征零}情形的方法来处理，而且还可使 $X_o$ 只有一个\Term{singular-point}{奇点}。下面就来说明这一方法。

\textbf{注记 12.7.2。}顺便指出，完成上述归约后，通常的论证还允许我们归约到以复\Term{number-field}{数域} $\mathbb C$ 为\Term{base-field}{基域}的情形，从而使用经典\Term{form}{形式}的 \Term{picard-lefschetz-formula}{Picard--Lefschetz 公式}；只是这个公式似乎还没有令人满意的文献可供引用。本研讨班对 \Term{picard-lefschetz-formula}{Picard--Lefschetz 公式}的证明，也将以类似的方式进行：先归约到\Term{transcendental}{超越}情形，再用真正的\Term{transcendental}{超越}方法处理。
\SourcePage{486}
因此，无论如何，这里给出的 12.5 的证明，进而 (10.4) 的证明，都具有\Term{transcendental}{超越}性质。另需指出，目前已知的 10.4 的另一种证明来自\Term{rigid-analytic}{刚性解析}理论 [24]；就此而言，它也应视为本质上“\Term{transcendental}{超越}的”。若能找到这些结果的纯\Term{algebra}{代数}证明，将十分有趣；但这是否可能，绝非显然！

\textbf{12.8。12.5 的证明：归约到 $S$ 为\Term{characteristic-zero}{特征零}的情形。}我们将使用一个论证，其原理与第 V 讲\SGAsourcegap{原 PDF 第486页此处在“V”后留空，缺少条款号}中“\Term{monodromy}{单值性}在 $\pi_1$ 上的\Term{action}{作用}本质上\Term{tame}{温和}”\Term{theorem}{定理}的证明相同。由 10.3，可以假设 $S$ \Term{complete}{完备}。设 $W$ 是\Term{residue-field}{剩余域}为 $k$ 的 Cohen $p$-环，并选取一个在\Term{residue-field}{剩余域}上诱导\Term{identity-map}{恒等映射}的\Term{homomorphism}{同态} $W\to V$（EGA $0_{\mathrm{IV}}$ 19.8.6）。令 $L=(x_\lambda)_{\lambda\in L}$ 为 $X_o$ 的\Term{singular-point}{奇点}\Term{set}{集合}，考虑 $X_o$ 在 $W$ 上的\Term{formal-moduli-scheme}{形式模概形} $\underline S$ 及\Term{special-fiber}{特殊纤维}为 $X_o$ 的“泛\Term{curve}{曲线}” $\underline f:\underline X\to\underline S$（第 VI 讲\SGAsourcegap{原 PDF 第486页“courbe universelle”所在句后的“VI”后留空，缺少条款号}）。在该处已看到，$\underline S$ \Term{isomorphism}{同构}于\Term{formal-power-series-ring}{形式幂级数环} $W[[T_1,\ldots,T_N]]$ 的\Term{spectrum}{谱}，并且可在 $\underline S$ 中找到一个\Term{normal-crossings-divisor}{正规交叉除子}
\[
 D=\sum_{\lambda\in L}D_\lambda,\tag{12.8.1}
\]
其\Term{irreducible-component}{不可约分支} $D_\lambda$ 由上述\Term{set}{集合} $L$ 编号，每个 $D_\lambda$ 由形如 $T_{i_\lambda}=0$ 的方程定义（若以整数集 $[1,r]$ 为 $L$ 编号，愿意的话可假定 $i_\lambda=\lambda$）。
\SourcePage{487}
$D_\lambda$ 是\Term{morphism}{态射} $\underline f:\underline X\to\underline S$ 的非\Term{smooth}{光滑}轨迹 $D^{\underline f}=\bigcup_{\lambda\in L}D_\lambda'$ 中经过 $X_o$ 的\Term{singular-point}{奇点} $x_\lambda$ 的\Term{irreducible-component}{分支} $D_\lambda'$ 经 $\underline f$ 的像。由构造，$f:X\to S$ \Term{isomorphism}{同构}于 $\underline f:\underline X\to\underline S$ 沿一个适当的 $W$-\Term{morphism}{态射} $g:S\to\underline S$ 的\Term{inverse-image}{逆像}：
\[
\begin{tikzcd}
 \underline X\arrow[d,"\underline f"'] & X\arrow[l]\arrow[d,"f"]\\
 \underline S & S\arrow[l,"g"]
\end{tikzcd}\tag{12.8.2}
\]
置
\[
 P=\underline{\operatorname{Pic}}_{\underline X/\underline S}^{0}.\tag{12.8.3}
\]
由于 $\underline f$ 的\Term{special-fiber}{特殊纤维}\Term{geometrically-reduced}{几何约化}且\Term{dimension}{维数}为 $1$，$\underline f$ 在\Term{dimension}{维数} $0$ 上\Term{cohomologically-flat}{上同调平坦}（EGA III 7.8.1），且 $P$ 在 $\underline S$ 上\Term{formally-smooth}{形式光滑}；后一点由 $H^2(X_o,\mathcal O_{X_o})=0$ 得出，因为 $X_o$ 的\Term{dimension}{维数}为 $1$。为简便起见，我们承认 $P$ 可由 $\underline S$ 上的\Term{group-scheme}{群概形}\Term{representable}{表示}，它必为具有\Term{connected}{连通}\Term{fibre}{纤维}的\Term{smooth}{光滑}\Term{group-scheme}{群概形}。这是 [22] 引用的 Mumford 一个未发表结果的特例；当底为\Term{trait}{离散赋值环谱}时，Raynaud 已在所引论文中大大推广了这个结果。由于 Mumford 和 Raynaud 的这些结果都还没有发表证明，读者也许更愿意使用 M. Artin 的\Term{generic}{一般}结果 [1]：只要起初的 $\underline f$ \Term{proper}{固有}、\Term{flat}{平坦}且在\Term{dimension}{维数} $0$ 上\Term{cohomologically-flat}{上同调平坦}，这些结果便给出 $P$ 可由 Artin 意义下的“\Term{algebraic-space}{代数空间}”\Term{representable}{表示}，
\SourcePage{488}
这对我们的目的已足够。置
\[
 U=\underline S-\operatorname{supp}D,\tag{12.8.4}
\]
于是 $U$ 是 $\underline S$ 中使 $\underline X$ 在其上\Term{smooth}{光滑}的最大\Term{open-subset}{开集}。因而
\[
 P|U=\underline{\operatorname{Pic}}_{\underline X_U/U}^{0}
\]
是 $U$ 上的\Term{abelian-scheme}{阿贝尔概形}，它沿 $g:\eta\to U$ 在 $\eta$ 上的\Term{inverse-image}{逆像}就是所要研究的 $\eta$ 上的\Term{abelian-scheme}{阿贝尔概形} $A_\eta$。更确切地说，在 $S$ 上有\Term{canonical}{典范}\Term{isomorphism}{同构} (12.1.10)：
\[
 P\times_{\underline S}S\xrightarrow\sim A^0.\tag{12.8.5}
\]
由于\Term{special-fiber}{特殊纤维} $P_o=A_o^0$ 对每个 $\ell$ 都是 $\ell$-\Term{divisible}{可除的}，$P$ 自身也如此，因此它的所有\Term{fibre}{纤维}都是\Term{abelian-variety}{阿贝尔簇}被\Term{torus}{环面}的\Term{extension}{扩张}。考虑 $\underline S$ 上的\Term{pro-group}{pro-群} $T_\ell(P)$；它沿 $g$ 的\Term{inverse-image}{逆像}正是 $T_\ell(A^0)$。应用 7.3 的构造，以一个“\Term{fixed-part}{固定部分}”和一个“\Term{toric-part}{环面部分}”定义这个\Term{pro-group}{pro-群}的\Term{filtration}{滤过}
\[
 T_\ell(P)\supset T_\ell(P)^f=T_\ell(\widehat P)
 \supset T_\ell(P)^t=T_\ell(T)\supset0.\tag{12.8.6}
\]
它沿 $g$ 的\Term{inverse-image}{逆像}是第 2、5 节中研究过的 $T_\ell(A^0)$ 的相应\Term{filtration}{滤过}。

考虑这个\Term{filtration}{滤过}在 $\underline S$ 的\Term{open-subset}{开集} $U$ 上的\Term{restriction}{限制}。$P_U$ 是 $U$ 上的\Term{smooth}{光滑}\Term{curve}{曲线}\Term{scheme}{概形}的\Term{jacobian}{Jacobian}，因此有\Term{canonical}{典范}\Term{self-duality}{自对偶性}，从而得到相应的\Term{pairing}{配对}
\SourcePage{489}
\[
 \varphi:T_\ell(P_U)\otimes T_\ell(P_U)\longrightarrow T_\ell(\mathbb G_{mU}),\tag{12.8.7}
\]
它是 $T_\ell(P_U)$ 的一个\Term{self-duality}{自对偶}，这里的\Term{self-duality}{自对偶}按\Term{barsotti-tate-group}{Barsotti--Tate 群}的意义理解。由 7.4.8 已知，对于这个\Term{pairing}{配对}，$T_\ell(P_U)^f$ 与 $T_\ell(P_U)^t$ 互为\Term{orthogonal}{正交补}。因此
\[
 E_U=T_\ell(P_U)/T_\ell(P_U)^t
\]
\Term{canonical}{典范}\Term{isomorphism}{同构}于 $T_\ell(P_U)^f=F_U$ 的\Term{dual}{对偶} Barsotti--Tate \Term{pro-group}{pro-群}，所以\Term{canonical}{典范}地延拓为 $\underline S$ 上的一个\Term{barsotti-tate-group}{Barsotti--Tate 群} $E$，即
\[
 E=D(F),\qquad F=T_\ell(P)^f.\tag{12.8.8}
\]
这时我们处于典型的\Term{blended-extension}{混合扩张}情形 (9.3)，可以重做 9.4 至 9.6 的构造；对于当时将这些构造写在底\Term{scheme}{概形}仅为一个\Term{trait}{离散赋值环谱}的陈旧而龌龊的背景下，我们深感抱歉！$T_\ell(P_U)$ 是 $E_U$ 被 $F_U$ 的一个\Term{blended-extension}{混合扩张}，其中 $E$、$F$ 是 $\underline S$ 上的\Term{barsotti-tate-group}{Barsotti--Tate 群}，它们分别是下列互为\Term{dual}{对偶}的\Term{extension}{扩张}：
\[
\begin{gathered}
\begin{tikzcd}[column sep=small]
0\arrow[r]&P\arrow[r]\arrow[d,equal]&F\arrow[r]\arrow[d,equal]&R\arrow[r]&0\\
&T_\ell(P)^t\simeq\operatorname{Hom}(M,T_\ell(\mathbb G_{m\underline S}))
&T_\ell(P)^f&&
\end{tikzcd}\\[1em]
\begin{tikzcd}[column sep=small]
0\arrow[r]&R\arrow[r]\arrow[d,equal]&E\arrow[r]&Q\arrow[r]\arrow[d,equal]&0\\
&D(R)&&M\otimes_{\mathbb Z}\mathbb Z_{\ell,\underline S}&
\end{tikzcd}
\end{gathered}\tag{12.8.9}
\]

\SourcePage{490}
（注意，由\Term{self-duality}{自对偶性} (12.8.7) 和上述\Term{orthogonality}{正交性}陈述，$R$ \Term{canonical}{典范}地\Term{self-dual}{自对偶}。）与所引各节的唯一区别在于，所得 $M_\ell\times M_\ell$ 上的\Term{pairing}{配对}取值于\Term{group}{群}
\[
 \bigl(\mathbb G_m(U)/\mathbb G_m(\underline S)\bigr)\otimes_{\mathbb Z}\mathbb Z_\ell;
\]
这个\Term{group}{群}不再\Term{canonical}{典范}\Term{isomorphism}{同构}于 $\mathbb Z_\ell$，而是借助\Term{isomorphism}{同构}
\[
 \mathbb Z^L\xrightarrow\sim\mathbb G_m(U)/\mathbb G_m(\underline S)\tag{12.8.10}
\]
\Term{canonical}{典范}\Term{isomorphism}{同构}于 $\mathbb Z_\ell^L$。这个\Term{isomorphism}{同构}将 $\mathbb Z^L$ 的标准基中下标为 $\lambda$ 的元素映为 $T_{i_\lambda}$ 在右边的类，其中 $T_{i_\lambda}$ 是\Term{divisor}{除子} $D_\lambda$ 的一个方程。于是得到\Term{canonical}{典范}\Term{pairing}{配对}
\[
 \underline u_\ell:M_\ell\otimes M_\ell\longrightarrow\mathbb Z_\ell^L.\tag{12.8.11}
\]
基本上不费力就能验证，所要研究的 $S$ 上的\Term{monodromy-pairing}{单值性配对} $u_\ell$ (12.4.7)，是 $\underline u_\ell$ 与\Term{homomorphism}{同态} $\mathbb Z_\ell^L\to\mathbb Z_\ell$ 的复合。后一个\Term{homomorphism}{同态}由下列\Term{homomorphism}{同态}在 $\mathbb Z$ 上与 $\mathbb Z_\ell$ 作\Term{tensor-product}{张量积}得到：
\[
\begin{tikzcd}[column sep=large]
 \mathbb G_m(U)/\mathbb G_m(\underline S)\arrow[r]\arrow[d,"\simeq"']
 &\mathbb G_m(\eta)/\mathbb G_m(S)\arrow[d,"\simeq"]\\
 \mathbb Z^L\arrow[r]&\mathbb Z.
\end{tikzcd}
\]
此\Term{homomorphism}{同态}由 $g:S\to\underline S$ 诱导，$g$ 将 $\eta$ 映入 $U$。它由一个元素
\SourcePage{491}
\[
 (d_\lambda)_{\lambda\in L}\in\mathbb Z^L,
 \qquad d_\lambda\geq1\quad(\lambda\in L)\tag{12.8.12}
\]
给出，其中每个 $d_\lambda$ 是 $\underline S$ 上的\Term{divisor}{除子} $D_\lambda$ 在 $S$ 中的\Term{inverse-image}{逆像}的\Term{multiplicity}{重数}；它其实只依赖于作为 $V$-\Term{algebra}{代数}的\Term{completion}{完备化} $\widehat{\mathcal O}_{X,x_\lambda}$（第 VI 讲\SGAsourcegap{原 PDF 第491页此处在第一次“VI”后留空，缺少条款号}）。还知道（第 VI 讲\SGAsourcegap{原 PDF 第491页此处在第二次“VI”后留空，缺少条款号}），对给定的 $\lambda$，$d_\lambda=1$ 当且仅当 $X$ 在点 $x_\lambda$ \Term{regular}{正则}。由于 12.5 假设了 $X$ \Term{regular}{正则}，即所有 $d_\lambda$ 都等于 $1$，所以 12.5 将由下列不再涉及 $X,S,g$ 的\Term{lemma}{引理}\Term{pushout}{推出}。

\begin{SGAresult}{\Term{lemma}{引理}}{12.8.13}
若 $\lambda\in L$，则 (12.8.11) 与 $\operatorname{pr}_\lambda$ 的复合等于\Term{bilinear-form}{双线性型} $\varphi_{x_\lambda}\otimes\varphi_{x_\lambda}$ (12.4.4)。
\end{SGAresult}
\begin{proof}
为证明这一陈述，考虑 $\underline S$ 在 $D_\lambda$ 的\Term{maximal-point}{极大点}的\Term{completed-strict-localization}{完备严格局部化} $\underline S^\lambda$，并令 $\underline X^\lambda=\underline X\times_{\underline S}\underline S^\lambda$。则\Term{special-fiber}{特殊纤维} $\underline X_o^\lambda$ 只有一个\Term{singular-point}{奇点} $y_\lambda$，仍为\Term{ordinary-quadratic-point}{普通二次点}。在相应的\Term{monodromy-pairing}{单值性配对}
\[
 u_\ell^\lambda:M_\ell^\lambda\otimes M_\ell^\lambda\longrightarrow\mathbb Z_\ell\tag{12.8.14}
\]
中，$M^\lambda$ 因而具有\Term{rank}{秩} $1$ 或 $0$。下面几个容易的事实，具体说明了上述构造在我们所需的特殊情形下对\Term{localization}{局部化}的行为，留给读者验证。

\textup{a)} $P\times_{\underline S}D_\lambda$ 中存在唯一的\Term{subtorus}{子环面} $T_\lambda$，使它在 $D_\lambda$ 的\Term{maximal-point}{极大点} $s_\lambda$ 上的\Term{restriction}{限制}为 $P_{s_\lambda}$ 的\Term{maximal-torus}{极大环面}，后者\Term{dimension}{维数}为 $0$ 或 $1$。

\SourcePage{492}
\textup{b)} 考虑\Term{inclusion}{包含}
$T_{\lambda o}\hookrightarrow T_o=P_o\text{ 的极大环面}$ 在\Term{character-group}{特征标群}上给出的\Term{transpose}{转置}\Term{homomorphism}{同态} $M\to M^\lambda$；再考虑仿照 12.4.3、以 $\underline X^\lambda$ 代替 $X$ 定义的型 $\varphi_{y_\lambda}:M^\lambda\to\mathbb Z$，此型确定到相差一个\Term{signature}{符号}。则下图在相差一个\Term{signature}{符号}的意义下\Term{commutative-diagram}{交换}：
\[
\begin{tikzcd}
 M\arrow[rr]\arrow[dr,"\varphi_{x_\lambda}"']&&M^\lambda\arrow[dl,"\varphi_{y_\lambda}"]\\
 &\mathbb Z&
\end{tikzcd}\tag{12.8.15}
\]
（事实上，在 $X_o$ 经过 $x_\lambda$ 的两个\Term{branch}{分支}的\Term{set}{集合}上选取一个次序，就可在 $\underline X_o^\lambda$ 的点 $y_\lambda$ 处作相应的选择；对于这样相互对应的两个选择，上图\Term{commutative-diagram}{交换}。）

\textup{c)} 下图\Term{commutative-diagram}{交换}：
\[
\begin{tikzcd}
 M_\ell\otimes M_\ell\arrow[rr]\arrow[dr,"\operatorname{pr}_\lambda\underline u_\ell"']
 &&M_\ell^\lambda\otimes M_\ell^\lambda\arrow[dl,"u_\ell^\lambda"]\\
 &\mathbb Z_\ell&
\end{tikzcd}
\]
其中记号如 12.8.13 和 (12.8.14) 所述。

由这些关系立即得出：为证明 12.8.13，只须对 $\underline X^\lambda$ 证明相应的关系
\[
 u_\ell^\lambda=(\varphi_{y_\lambda}\otimes\varphi_{y_\lambda})
 \otimes_{\mathbb Z}\mathbb Z_\ell,
\]
\SourcePage{493}
而这就是 $\underline S^\lambda$ 上的 $\underline X^\lambda$ 的情形下的 12.5。但 $\underline S$ 是一个\Term{characteristic-zero}{特征零}的\Term{trait}{离散赋值环谱}。所述归约至此完成。
\end{proof}

\textbf{注记 12.9。}注意，这一归约还同时将我们归约到了 $X_o$ 只有一个\Term{singular-point}{奇点}的情形，这正是经典 \Term{picard-lefschetz-formula}{Picard--Lefschetz 公式}所考虑的情形。此外，这时 $X_o$ 有一个或两个\Term{irreducible-component}{不可约分支}，分别对应于 $\operatorname{rank}M=1$ 和 $\operatorname{rank}M=0$；后一种正是\Term{generic-fiber}{泛纤维}的\Term{jacobian}{Jacobian}具有\Term{good-reduction}{良约化}的情形。

\textbf{12.10。}12.8 中的证明还给出一个比 12.5 更\Term{generic}{一般}的结果，无须假设 $X$ \Term{regular}{正则}：引入 (12.8.12) 中的整数 $d_\lambda\geq1$，则所有\Term{monodromy-pairing}{单值性配对}均由\Term{pairing}{配对}
\[
 u=\sum_{\lambda\in L}d_\lambda\,
 \varphi_{x_\lambda}\otimes\varphi_{x_\lambda}:
 M\otimes M\longrightarrow\mathbb Z\tag{12.10.1}
\]
\Term{derive}{导出}。
